\chapter{An exact evaluator and a trustworthy quote}\label{ch:evaluator}
An evaluator's task is narrow: given a declared baseline and accepted finite
action, return the corresponding value without moving the model. A narrow
task is easier to audit because unexpected side effects have no legitimate
place to hide.

\section{Independent arithmetic routes}
For the Gaussian domain, two direct routes should agree:
\begin{align}
 E_G&=V(z)-V(z+\delta_G),\\
 E_G&=-\mu(z)^T\delta_G-\tfrac12\delta_G^TH\delta_G.
\end{align}
Local affected-factor evaluation provides a third comparison with the full
diagnostic difference. These routes share a mathematical specification but
can expose different implementation mistakes: wrong signs, omitted curvature,
incomplete support or a stale baseline.

A test that compares a function with another call to the same function is
not an independent check. A fixed rational example with an independently
written expected value is often more revealing than thousands of repeated
calls through one mistaken helper.

\section{Exact algebra and numerical representation}
The formula is exact mathematics. Floating-point evaluation is not exact
real arithmetic. A future protocol must specify whether representation is
rational, fixed-point or floating-point, and what refusal or tolerance rules
apply near zero, affordability and conservation boundaries.

Reporting a residual of zero on selected fixtures does not prove all real
inputs will produce exact machine zero. Conversely, choosing a convenient
tolerance after observing a failure can weaken the scientific contract.
The numerical policy must be fixed before outcomes can influence it.

\section{Value, attribution and settlement have different outputs}
The group value is a single field difference. A common-path evaluator can
add $R_a$ for each child and verify $\sum_aR_a=E_G$. A settlement service
would additionally require an accepted owner/policy map and a valid event
identity. Combining all three in one return value called reward invites
the caller to assume an authority that the evaluator never established.

The preserved candidate code already distinguishes physical constant-rate
path semantics from endpoint interpolation. That distinction is reusable
as an idea. Its mere presence in a candidate module does not authorize
spendable Gaussian capacity or prove every interpretation guard sufficient.

\section{A quote must describe what was accepted}
For a one-dimensional accepted schedule, the record can identify the full
function on its accepted domain. A measured shortfall can select a point
on that precommitted function. For a simultaneous group, changing one child
may alter every path attribution. The applicable multidimensional or
re-epoching contract must therefore be declared rather than extrapolated
from the simplest single-action schedule.

An overshoot outside a committed domain is a protocol violation. It does
not justify an improvised positive credit, nor an arbitrary punishment
presented as physical EBU. Physical measurement, account calculation and
sanctions remain separate questions.

\section*{Worked evaluator audit}
At $(18,2)$, reference $(10,10)$, scales $(2,2)$ and quantity four,
the endpoint is $(14,6)$. The potential route gives $16-4=12$.
The tangent is sixteen and the curvature correction four, again giving
twelve. A returned value of sixteen suggests the finite evaluator has
silently become a first-order approximation. A returned value of fourteen
for two two-unit children suggests independent singleton values were added.

Neither error is cured by increasing decimal precision. The wrong
mathematical object was computed.

\section*{Chapter recap}
Exact evaluation needs a complete baseline, support, increment and numerical
policy. A successful calculation is not permission to execute, and a closing
decomposition is not authorization to credit an owner.
